\section{Scalar outputs of sparse conic and linear programs}\label{sec:realizations}

The matrix lower bounds of the preceding sections need not disappear
when the scalar must be an ordinary optimization output: we construct sparse SOCPs
and LPs in which a designated optimizer coordinate, the optimal value, or
a Newton decrement encodes the same hidden Forrelation amplitude.  We also
determine how the bounds change when the output is a sample.  The clock construction builds on the Forrelation
matrix-function reductions of Montanaro and Shao \citep{MontanaroShao2023}
and the inverse-sampling separation of Gr\o nlund and Larsen
\citep{GronlundLarsen2025}; neither is new here.

\subsection{A cyclic inverse with public norms}

Let $q=2^r$, $r\geq1$, and let $n=r\ell$ be the number of work qubits.
A two-function Forrelation circuit $V$ has three global Walsh--Hadamard
transforms and two hidden diagonal sign gates.  Its real amplitude
$\Phi=\langle0^n|V|0^n\rangle$ obeys the promise
\[
 |\Phi|\leq1/100\quad\text{or}\quad\Phi\geq3/5.
\]
The randomized query lower bound is
$\Omega(2^{n/2}/n)$ \citep{AaronsonAmbainis2018}.
Split each global transform into $\ell$ transforms on disjoint $r$-qubit
blocks.  Padding with at most one identity layer makes the number $T$ of
forward layers even, where $T=3\ell+2$ or $3\ell+3$.  Each layer is $q$-sparse.

Use a clock of length $L=3T$.  The transitions run $V$ forward, hold its
final state for $T$ steps, then run its gates in reverse transposed order.
Let $U$ be this real orthogonal clock transition.  Let $W_t$ be the
product of the transitions up to clock time $t$, with $W_0=I$.  The
orthogonal gauge
$W=\sum_{t=0}^{L-1}|t\rangle\langle t|\otimes W_t$ satisfies
\begin{equation}\label{eq:history-gauge}
 W^TUW=S\otimes I,\qquad U^L=I,
\end{equation}
where $S|t\rangle=|t+1\bmod L\rangle$.  In particular $W_t=V$ for
$T\leq t<2T$.  For $K\geq3$, define
\begin{equation}\label{eq:cyclic-M}
 \gamma=\frac{K-1}{K+1},\qquad
 \alpha_K=\log\frac{K+1}{K-1},\qquad
 M=\frac{I-\gamma U}{1+\gamma},\qquad e=|0,0^n\rangle.
\end{equation}
Here $M$ need not be symmetric; the associated Newton Hessian will be SPD.
Since $L$ is even, the spectrum of $S$ contains $1$ and
$-1$.  Thus
\begin{equation}\label{eq:cyclic-spectrum}
 \norm M=1,\qquad \sigma_{\min}(M)=K^{-1},\qquad
 \kappa_2(M)=K,\qquad \nnz(M_{i,:}),\nnz(M_{:,j})\leq q+1.
\end{equation}

\begin{lemma}[Inverse history and full sampling access]\label{lem:cyclic-history}
Write $z_0=\gamma^{3T}$, $S_0=\sum_{t=0}^{3T-1}\gamma^{2t}$,
$S_I=\sum_{t=T}^{2T-1}\gamma^{2t}$, and $z=\gamma^{2T}$.  Then
\begin{align}
 x_*:=M^{-1}e
 &=\frac{1+\gamma}{1-z_0}\sum_{t=0}^{3T-1}\gamma^tU^te,
 & R^2:=\norm{x_*}^2&=K\frac{1+z_0}{1-z_0},\label{eq:history}\\
 p:=S_I/S_0&=\frac{z}{1+z+z^2},
 & K_{\rm eff}:=K/R&=\sqrt{K\frac{1-z_0}{1+z_0}}.
 \label{eq:plateau}
\end{align}
Each call to a full SQ oracle for $M$, for
\begin{equation}\label{eq:cyclic-hessian}
 H_0=2M^TM=\frac{2}{(1+\gamma)^2}
       [(1+\gamma^2)I-\gamma(U+U^T)],
\end{equation}
or for $M^Te$ can be simulated with at most a constant number of hidden
sign queries.  The Hessian has condition number $K^2$ and sparsity
at most $2q+1$.
\end{lemma}

\begin{proof}
Multiplying the finite geometric series by $I-\gamma U$ proves the
inverse formula.  Different clock positions are orthogonal and each
$U^te$ is a unit vector.  Hence
$R^2=(1+\gamma)^2 S_0/(1-z_0)^2$; summing the geometric progression and
using $(1+\gamma)/(1-\gamma)=K$ gives \eqref{eq:history}.
Partitioning $S_0$ into three equal clock segments gives \eqref{eq:plateau}.
Normality of $U$ gives \eqref{eq:cyclic-hessian} and the spectrum assertion.

For the access claim, the diagonal, forward-clock, and reverse-clock
supports are disjoint.  A transition is either a public Hadamard block,
a hidden diagonal sign, or an identity.  Each of its rows and columns has
squared norm one, and its entry magnitudes and locations are public.  Therefore every row
and column of $M$ has squared norm $(1+\gamma^2)/(1+\gamma)^2$; every
row and column of $H_0$ has the same public squared norm
$4[(1+\gamma^2)^2+2\gamma^2]/(1+\gamma)^4$.
Their row-norm sampling distributions are uniform.  A within-row sample
chooses a clock block with public probabilities, then samples a Hadamard
row uniformly or returns the single location of a sign entry.
The corresponding column rules are identical.  Locations and norms need
no hidden-input query; a value uses at most one.  The vector $M^Te$ is a
public-norm mixture of the clock-zero basis vector and a row of the
last transition.  The same simulation supplies its SQ interface.
\end{proof}

Fix a sufficiently small constant $\delta_0>0$.  For
$0<\delta\leq\delta_0$, choose an allowed $T$ nearest to
$\log(1/\delta)/\alpha_K$.  Allowed lengths have bounded gaps and
$\gamma\geq1/2$, so, with universal constants,
\begin{equation}\label{eq:scalar-clock-parameters}
 p=\Theta(\delta^2),\quad R=\Theta(\sqrt K),\quad
 \ell=\frac{\log(1/\delta)}{3\alpha_K}+O(1),\quad
 N=3Tq^\ell.
\end{equation}
Taking $\delta_0$ small ensures $\ell\geq1$.  All lower bounds below
hold for the dimensions $N$ selected in this way, not for every
prescribed $N$.

\subsection{An optimizer coordinate and approximately optimal feasible points}

Let $Q_{N+1}=\{(t,u):t\geq\norm u\}$ be the Lorentz cone.  Consider
\begin{equation}\label{eq:one-cone-program}
 \max_{x,u}\ 2e^Tu
 \quad\text{subject to}\quad u=Mx,\quad(1,u)\in Q_{N+1}.
\end{equation}
Its unique optimizer is $(x_*,e)$, and its optimal value is two for every
hidden input, so estimating this value is trivial.  Substituting $u=Mx$ gives the barrier
$\phi(x)=-\log(1-\norm{Mx}^2)$.  This barrier has parameter one, whereas
the standard barrier of the full Lorentz cone has parameter two.  It is self-concordant as an affine restriction.  In the coordinates $u=Mx$, the
squared dual local norm of its gradient is
$2\norm u^2/(1+\norm u^2)\leq1$, proving the parameter bound.
The analytic center is zero for every hidden input.  At this center, for
$F_\eta(x)=\phi(x)-2\eta e^TMx$,
\begin{equation}\label{eq:first-newton}
 \nabla F_\eta(0)=-2\eta M^Te,\qquad
 \nabla^2 F_\eta(0)=H_0,\qquad
 d_\eta=\eta M^{-1}e,\qquad \Lambda_\eta^2=2\eta^2.
\end{equation}
For $\eta<1$, this full Newton step is strictly feasible; for
$\eta\leq1/(4\sqrt2)$ its decrement is at most $1/4$.
The central path is
$x(\eta)=\rho(\eta)x_*$, where
$\rho/(1-\rho^2)=\eta$ and $0<\rho<1$.
These identities concern the reduced Hessian, not an augmented KKT
matrix with redundant free variables.

For $t$ in the hold segment, or plateau, $I:=\{T,\ldots,2T-1\}$, where
$W_t=V$, set $i_t=(t,0^n)$ and
\begin{equation}\label{eq:public-readout}
 a=\sum_{t\in I}\frac{\gamma^t}{R\sqrt{S_I}}e_{i_t},
 \qquad h=\sqrt p.
\end{equation}
Then $\norm a=1/R$ and $a^Tx_*=h\Phi$.
To make this linear observable an optimizer coordinate, we append free
variables satisfying
\begin{equation}\label{eq:accumulator}
 y_T=a_{i_T}x_{i_T},\qquad
 y_{t+1}=y_t+a_{i_{t+1}}x_{i_{t+1}},\qquad
 w=y_{2T-1}.
\end{equation}
The chain has $O(T)$ rows of at most three entries.  Each
selected column of $x$ gains one nonzero entry.  These equalities and
their norms are public, so full SQ access to the augmented formulation
mixes public data with the interface of
Lemma~\ref{lem:cyclic-history}.  The maximum row and column sparsity
remains $s=\Theta(q)$.  The free accumulator variables add no cone and can
be eliminated when computing barrier derivatives.

\begin{theorem}[Sparse scalar coordinate separation]\label{thm:scalar-coordinate}
For the parameters in \eqref{eq:scalar-clock-parameters}, estimating the
optimizer coordinate $w_*=h\Phi$ to additive error $c\delta$, for a
sufficiently small universal $c$, requires
\begin{equation}\label{eq:scalar-coordinate-lower}
 Q_{\rm classical}=\Omega\left(\frac{q^{\ell/2}}{r\ell}\right)
\end{equation}
full-SQ input queries at bounded error.  The same bound applies to
returning this coordinate of any exactly feasible point with objective gap
$\tau\leq c_0\delta^2/K$, for a sufficiently small constant $c_0$.
A coherent input oracle is \emph{source-equivalent} if each of its calls
can be simulated with $O(1)$ hidden Forrelation sign queries and vice
versa.  With such an oracle, estimating the exact optimizer
coordinate to the stated accuracy takes $O(1)$ hidden
queries and $\operatorname{polylog}N$ public gates.
\end{theorem}

\begin{proof}
The coordinate identity follows by taking the overlap of
\eqref{eq:history} with \eqref{eq:public-readout}.  Its low and high
intervals are separated by $(3/5-1/100)h=\Theta(\delta)$.
Each formulation query costs a constant number
of hidden sign queries, so the
Forrelation lower bound proves \eqref{eq:scalar-coordinate-lower}.
If $(\widehat x,\widehat u)$ is exactly feasible and has gap $\tau$, then
\[
 \norm{\widehat u-e}^2\leq2-2e^T\widehat u\leq\tau,
 \qquad
 |a^T\widehat x-h\Phi|
 \leq\frac{K}{R}\sqrt\tau\leq\sqrt{K\tau}.
\]
Exact accumulator feasibility equates the returned coordinate with
$a^T\widehat x$.  The specified $\tau$ therefore preserves the promise
gap.  Points that violate the
equalities or the cone constraint by a small residual would need a
separate error analysis.

A coherent query to a designated transition entry recovers any desired
hidden sign, and conversely each coherent matrix-oracle call uses a
constant number of sign queries.  Run the source circuit and estimate
its real amplitude to a sufficiently small constant error, then multiply
by public $h$.
\end{proof}

The feasible-point lower bound uses such a point only as a source of one
accurate coordinate, which the quantum algorithm of
Theorem~\ref{thm:scalar-coordinate} estimates without outputting the vector.

Since $\alpha_K=\Theta(1/K)$, the logarithm of
\eqref{eq:scalar-coordinate-lower} is
$\Omega(K\log s\log(1/\delta))$, even with the denominator $r\ell$ and
the $O(1)$ rounding term in $\ell$ kept.  This exponent
grows with $\log s$, not with the ambient dimension.  With
$k$-Forrelation, which uses a fixed number $k$ of hidden sign tables, the
lower bound becomes $\widetilde\Omega_k(N^{1-1/k})$, a larger power of $N$
for $k\geq3$ (Corollary~\ref{cor:fixed-k}).

\begin{corollary}[Fixed-order Forrelation]\label{cor:fixed-k}
Fix $k\geq2$.  Replace $V$ by the $k$-Forrelation circuit and use the
primary promise
\[
 |\Phi|\leq\alpha_k:=2^{-5k-1}
 \quad\text{or}\quad \Phi\geq\beta_k:=2^{-5k}.
\]
Then $T=(k+1)\ell+O(k)$ and
\[
 \ell=\frac{\log(1/\delta)}{(k+1)\alpha_K}+O_k(1).
\]
At coordinate accuracy $c(\beta_k-\alpha_k)\delta$, the lower bound is
\begin{equation}\label{eq:fixed-k-lower}
 \Omega_k\left(\frac{q^\ell}{r\ell}\right)^{1-1/k}
 =\widetilde\Omega_k(N^{1-1/k}).
\end{equation}
For exactly feasible points, an objective gap of $c_k\delta^2/K$, with
$c_k=2^{-O(k)}$, suffices for the same decision.
The coherent scalar algorithm uses $2^{O(k)}$ hidden-input queries.
\end{corollary}

\begin{proof}
Theorem 1.3 and Corollary 1.4 of Bansal and Sinha
\citep{BansalSinha2021} give, for fixed $k$, the stated randomized
lower bound and a $2^{O(k)}$-query bounded-error algorithm.  A problem whose
answer depends only on $|\Phi|$ may also allow the high case
$\Phi\leq-\beta_k$; this only enlarges the promise domain, so the lower
bound still applies.  The $k+1$ global Hadamards give
the claimed clock length.  All remaining identities and oracle
simulations are unchanged.  Estimating the amplitude to error
$O(\beta_k-\alpha_k)$ proves the numerical upper bound, and the
feasible-point error $\sqrt{K\tau}$ proves the last accuracy claim.
\end{proof}

\subsection{A hard optimal value from a tilt with public norm}

The optimal value of \eqref{eq:one-cone-program} is public, but adding a
public tilt to the objective makes the value depend on $\Phi$.  Define
\[
 g=M^{-T}a,\qquad G=\norm g.
\]
Although $g$ depends on the hidden matrix $M$, the plateau $I$ makes
$G$ exactly public (Lemma~\ref{lem:tilt-norm}), so the quadratic tilt
term in the optimal value is known rather than an unknown error.

\begin{lemma}[Public tilt norm and plateau optimality]\label{lem:tilt-norm}
Let $m=(I-\gamma S)/(1+\gamma)$ and
$\bar a=\sum_{t\in I}\gamma^t|t\rangle/(R\sqrt{S_I})$.
Then
\begin{equation}\label{eq:public-G}
 G=\norm{m^{-T}\bar a}=\Theta(\sqrt K),\qquad
 h\leq G\leq K/R\leq\sqrt K.
\end{equation}
For any clock vector $\bar c$ supported on $I$, with
$b_c=m^{-T}\bar c\ne0$,
\begin{equation}\label{eq:plateau-optimality}
 \frac{|(b_c)_0|}{\norm{b_c}}
 \leq\left(\sum_{t=0}^{T-1}\gamma^{-2t}\right)^{-1/2}
 =\Theta(\delta/\sqrt K).
\end{equation}
Thus, up to constant factors, the readout \eqref{eq:public-readout}
maximizes the normalized inverse-transpose overlap
\eqref{eq:plateau-optimality} among clock readouts supported on the
plateau $I$.
\end{lemma}

\begin{proof}
On the plateau $W_t=V$.  Applying \eqref{eq:history-gauge} gives
\[
 W^Ta=\bar a\otimes V^T|0^n\rangle,\qquad
 W^Tg=(m^{-T}\bar a)\otimes V^T|0^n\rangle.
\]
Hence $G$ does not depend on the hidden input.  Its value
follows from
the finite geometric inverse of $m^T$, requiring no hidden-input queries.
If $b=m^{-T}\bar a$, then
$b_0=\bar a^Tm^{-1}|0\rangle=\sqrt{S_I/S_0}=h$.
Since $\bar a_t=0$ for $0\leq t<T$, the equation
$b_t-\gamma b_{t+1}=0$ gives $b_t=h\gamma^{-t}$ there.  Hence
\begin{equation}\label{eq:G-lower}
 G^2\geq h^2\sum_{t=0}^{T-1}\gamma^{-2t}
 =\frac{1-z}{1+z+z^2}\frac{\gamma^2}{1-\gamma^2}
 =\frac{1-z}{1+z+z^2}\frac{(K-1)^2}{4K}.
\end{equation}
For the chosen $z=\Theta(\delta^2)$ this is $\Omega(K)$.
The upper bound is $\norm{m^{-1}}\norm{\bar a}=K/R$.
The same homogeneous recurrence for $b_c$ proves
\eqref{eq:plateau-optimality}; evaluating its geometric sum gives the
last order estimate.  For any public readout on $(t,0^n)$ over this
plateau, its hidden overlap is $(b_c)_0\Phi$ and its direction norm is
$\norm{b_c}$.  Our choice has ratio $h/G=\Theta(\delta/\sqrt K)$.
\end{proof}

\begin{theorem}[Optimal value and direct Newton decrement]\label{thm:tilted-value}
Keep the accumulator equalities \eqref{eq:accumulator} and replace the
objective in \eqref{eq:one-cone-program} by $e^Tu+\lambda w$, where
$\lambda=1/(2G)$.
Then
\begin{equation}\label{eq:tilted-value}
 \OPT^2=\frac54+\frac hG\Phi,\qquad \frac12\leq\OPT\leq\frac32.
\end{equation}
Estimating the optimal value to additive error
$c(\beta-\alpha)\delta/\sqrt K$ has lower bound
\eqref{eq:scalar-coordinate-lower} for the two-function family and
\eqref{eq:fixed-k-lower} for fixed-order Forrelation, where $\alpha$ and
$\beta$ are the corresponding source thresholds.
With the source-equivalent coherent oracle, an algorithm attains that
accuracy with $O(1)$ or $2^{O(k)}$ hidden-input queries, respectively.

At the public analytic center, with objective multiplier $\eta>0$, the
squared Newton decrement is
\begin{equation}\label{eq:tilted-decrement}
 \Lambda_\eta^2
 =\frac{\eta^2}{2}\left(\frac54+\frac hG\Phi\right).
\end{equation}
The same lower bounds hold for estimating $\Lambda_\eta^2$ to additive
error $c(\beta-\alpha)\eta^2\delta/\sqrt K$, even when the algorithm is
given full SQ access to the reduced Hessian and its right-hand side.
\end{theorem}

\begin{proof}
Eliminating the equalities gives $w=g^Tu$.  Maximizing
$(e+\lambda g)^Tu$ on the unit ball gives $\OPT=\norm{e+\lambda g}$.
Since $e^Tg=h\Phi$ and $\lambda^2G^2=1/4$, this proves
\eqref{eq:tilted-value}.  The bound $h/G\leq1$ proves its value range.
The low and high cases have respectively absolute squared deviations
from $5/4$ at most $\alpha h/G$ and at least $\beta h/G$.
On $[1/2,3/2]$ squaring and taking square roots have constant Lipschitz
bounds, proving the accuracy transfer.  The objective has only the
public $e$ and $w$ coefficients, so its access adds no hidden-input query.
The coherent algorithm outputs
$\sqrt{5/4+(h/G)\widehat\Phi}$ after estimating $\Phi$ to a small
multiple of $\beta-\alpha$; clipping $\widehat\Phi$ to $[-1,1]$ is harmless.

The reduced gradient is $-\eta c$, where
$c=M^Te+\lambda a=M^T(e+\lambda g)$.
Multiplying by $(2M^TM)^{-1}$ gives
\eqref{eq:tilted-decrement}.  The supports of $M^Te$ and $a$ occupy
disjoint public clock regions, and both component norms are public.
Thus $\SQ(c)$ is a known mixture, with at most one hidden-input query for a
value.  Lemma~\ref{lem:cyclic-history} supplies full SQ access to the
Hessian.  This proves the last claim.
\end{proof}

The eliminated tilt coefficients stay bounded: every singular value of
$m^{-T}$ is at least one, so
$G\geq\norm{\bar a}=1/R$ and hence $\norm{\lambda a}\leq1/2$.
The optimality statement in Lemma~\ref{lem:tilt-norm} covers only readouts
on this plateau with a normalized objective direction of constant scale,
and does not rule out other scalar reductions.
For comparison, if only the bound $G^2\leq K$ were used, one could take
the smaller public tilt $\lambda=\theta h/K$.  Then
\[
 \OPT_\lambda^2=1+2\theta h^2\Phi/K+\theta^2h^2G^2/K^2.
\]
Taking $\theta\leq c(\beta-\alpha)$ separates the cases at accuracy
$\Theta(\theta(\beta-\alpha)\delta^2/K)$, even without computing $G$.
Using the exact public $G$ therefore increases the gap between the low-
and high-case optimal values from order $\delta^2/K$ to $\delta/\sqrt K$.

For requested value accuracy $\varepsilon$ in the two-function case,
choose $\delta=\Theta(\varepsilon\sqrt K)$ within
the small-$\delta$ regime.  The lower bound then grows as
\[
 \exp\left(\Omega\left(K\log s\,
       \log\frac1{\varepsilon\sqrt K}\right)\right),
\]
and \eqref{eq:scalar-coordinate-lower} gives the exact denominator.
For the reduced-Hessian condition number $\kappa_H=K^2$, substitute
$K=\sqrt{\kappa_H}$.  This realizes the lower bound as an optimal value
in a high-accuracy regime; it does not replace
Theorem~\ref{thm:clock-lower}, which holds over a wider parameter range.

\subsection{Two LP realizations and their access requirements}

The decrement \eqref{eq:tilted-decrement} also arises in an LP with box
constraints.  Set
$c=M^Te+\lambda a$ as above and consider
\[
 \max_x c^Tx\quad\text{subject to}\quad-\mathbf1\leq Mx\leq\mathbf1.
\]
The logarithmic barrier of the paired constraints,
$\phi_{\rm box}(x)=-\sum_i\log(1-(Mx)_i^2)$, has analytic center zero
and Hessian $H_0$ there.  Its decrement is exactly
\eqref{eq:tilted-decrement}.  Full SQ access to $[M;-M]$ and $c$ is simulated as before,
so the decrement lower bound under direct full SQ access to the Newton
system already holds for an LP with scalar inequality constraints.  We do not claim that
the optimal value of the box LP equals that of the SOCP.

For the unperturbed objective $2e^TMx$, the box barrier subproblem in
$u=Mx$ is
\[
 F_\eta(u)=-\sum_i\log(1-u_i^2)-2\eta e^Tu.
\]
Separability gives $u(\eta)=\rho(\eta)e$, with
$\rho/(1-\rho^2)=\eta$.  Therefore the central path lies on the same ray as the Newton step in
\eqref{eq:first-newton}, and its
accumulator coordinate is $\rho h\Phi$.  Its Hessian is at least $2I$ in
$u$ coordinates.  Consequently, if $F_\eta(\widehat u)$ exceeds its minimum by at most
$\tau$, then
$\norm{\widehat u-\rho e}\leq\sqrt\tau$ and the accumulator error is
at most $\sqrt{K\tau}$.  For any fixed public $\eta>0$, the coordinate lower bound
therefore still holds for points with barrier-subproblem gap at most
$c\rho(\eta)^2(\beta-\alpha)^2\delta^2/K$, for a sufficiently small
absolute constant $c>0$.  Again, the accumulator equalities must hold
exactly, or the reported scalar must be computed as $a^T\widehat x$.
The box barrier has parameter $N$, since each paired interval barrier
has parameter one; this realization gives no small iteration bound.

A second LP replaces the full-dimensional box by a one-dimensional
affine slice.  Let $\widetilde a=Ra$, so $\norm{\widetilde a}=1$, and set
\begin{equation}\label{eq:affine-lp}
 \max_{x,t}\ \widetilde a^Tx
 \quad\text{subject to}\quad Mx=te,\quad-1\leq t\leq1.
\end{equation}
Replace $a$ by $\widetilde a=Ra$ in the accumulator equalities
\eqref{eq:accumulator} and use $w$ as the objective.  This rescales only
the readout coefficients of $x$, makes the objective one-sparse, and
leaves the projection of the feasible set onto $(x,t)$ unchanged.

\begin{theorem}[Affine-slice value and matched upper exponent]\label{thm:affine-lp}
The LP \eqref{eq:affine-lp} has a unit-norm objective vector and
full-SQ input access with a constant number of hidden-input queries per call.  Its
exact value is
\begin{equation}\label{eq:affine-value}
 \OPT=Rh|\Phi|=\Theta(\sqrt K\,\delta)|\Phi|.
\end{equation}
Its equality-constraint matrix without the accumulator rows,
$A_{\rm eq}=[M\ -e]$, has
$K\leq\kappa_2(A_{\rm eq})\leq\sqrt2K$.
Estimating the value to additive error $c(\beta-\alpha)\sqrt K\,\delta$
has the classical lower bounds \eqref{eq:scalar-coordinate-lower} and
\eqref{eq:fixed-k-lower}, and coherent algorithms achieve it with the
corresponding constant or $2^{O(k)}$ number of hidden-input queries.  In
particular, for fixed $k$ and constant additive value accuracy, one can
take $\delta=\Theta_k(K^{-1/2})$, which gives a lower bound with exponent
$\Omega_k(K\log s\log K)$.

For the two-function family, at additive accuracy $\varepsilon$ with
$\varepsilon/\sqrt K$ sufficiently small, the logarithms of the classical
lower and upper bounds agree up to constant factors:
\begin{equation}\label{eq:affine-matched}
 \log Q=\Theta\left(K\log s\log\frac{\sqrt K}{\varepsilon}\right)
\end{equation}
when the displayed exponent dominates the polynomial and rounding
factors.  The exact lower bound remains
\eqref{eq:scalar-coordinate-lower}, with
$\ell=[\log(\sqrt K/\varepsilon)+O(1)]/(3\alpha_K)+O(1)$.
The bracketed $O(1)$ term comes from constant accuracy factors and
changes $\ell$ by $O(K)$; the final $O(1)$ term accounts for rounding to
an allowed length.
\end{theorem}

\begin{proof}
Invertibility forces $x=tM^{-1}e$, so maximizing over $|t|\leq1$
gives \eqref{eq:affine-value}.  The added column $-e$ changes only one
public row norm of $M$; its location, norm, and all conditional mixture
weights are public.  The objective $\widetilde a$ is also public.
Moreover,
\[
 K^{-2}I\preceq A_{\rm eq}A_{\rm eq}^T=MM^T+ee^T\preceq2I.
\]
The $K^{-2}$ eigenspace of $MM^T$ has multiplicity $q^\ell\geq2$.
A vector in it orthogonal to $e$ preserves the smallest eigenvalue.
The largest singular value is at least $\norm M=1$, proving the
condition bounds.  No such bound is asserted for the accumulator-augmented
KKT matrix.  The value identity and the source gap prove the lower
and coherent upper bounds exactly as before.

For the classical upper bound, apply Theorem~\ref{thm:general-overlap}
to $(M,\widetilde a,e)$, using the exact public
$R_e=R=\Theta(\sqrt K)$.  The theorem permits
$0<\varepsilon\leq R/2$, which contains the chosen small-$\delta$
value-accuracy regime even when $\varepsilon>1/2$.
The absolute-value function is 1-Lipschitz,
so the theorem estimates \eqref{eq:affine-value} at cost
\[
 O\left(R^2\varepsilon^{-2}\log(1/\zeta)\,
       (q+2)^{O(K\log(2R/\varepsilon))}\right).
\]
This yields the upper exponent in \eqref{eq:affine-matched} without
constructing a nullspace basis, an optimizer, or an SQ interface to
$M^{-1}e$.  Substituting $\delta=\Theta(\varepsilon/\sqrt K)$ in
the lower family yields the matching exponent.
\end{proof}

In its intrinsic coordinate $t$, the affine slice has the barrier
$\varphi(t)=-\log(1-t^2)$ with parameter one.  Indeed,
\[
 \varphi'=\frac{2t}{1-t^2},\quad
 \varphi''=\frac{2(1+t^2)}{(1-t^2)^2},\quad
 \varphi'''=\frac{4t(3+t^2)}{(1-t^2)^3}.
\]
The gradient bound is $(\varphi')^2/\varphi''=2t^2/(1+t^2)\leq1$.
Squaring the self-concordance inequality
$|\varphi'''|\leq2(\varphi'')^{3/2}$ reduces to
$(1-t^2)^2(t^2+2)\geq0$.
At the analytic center, the tangent space is
$\operatorname{span}\{(M^{-1}e,1)\}$, the barrier Hessian in $t$ is two,
and the projected objective is $Rh\Phi$.  Therefore
\begin{equation}\label{eq:affine-decrement}
 \Lambda_\eta^2=\frac{\eta^2}{2}R^2h^2\Phi^2
              =\frac{\eta^2}{2}\OPT^2.
\end{equation}
Estimating this squared decrement to additive error
$c(\beta^2-\alpha^2)\eta^2 K\delta^2$ has the same source lower bound.
Here the scalar projected right-hand side is \emph{not} given, since it
already contains the hard quantity.  If a nullspace
basis and the projected objective were supplied, the program would reduce
to an explicit problem on an interval.  Thus, in this parameter-one realization, the hardness lies in eliminating the equality
constraints and reading out the scalar, not in repeated high-dimensional
barrier computations.

\subsection{Replacing the large cone by three-dimensional cones}

A binary tree of three-dimensional Lorentz cones can replace the single
large cone.  Let $D$ be the smallest power of two at
least $N$, pad $u$ by zeros to $D$ leaves, and for each internal node
$v$ impose $(t_v,z_{v,0},z_{v,1})\in Q_3$, where its children are leaf
coordinates or lower node variables.  Fix the root $t$ to one.
Recursive use of $t_v^2\geq z_{v,0}^2+z_{v,1}^2$ shows that every feasible
$u$ satisfies $\norm u\leq1$.  Choosing each non-root $t_v$
as the norm of the leaf coordinates below $v$ proves the converse.  Public two-sparse equalities that copy
shared variables put the constraints in standard product-cone form.  The
construction uses $D-1$ cones, each new variable appears in a constant
number of constraints, and all new SQ data are public, so all
optimizer-coordinate and optimal-value results above still hold.

The first Newton system can also be computed exactly.  After the copied
variables are identified, the product barrier is
\[
 \Psi(t,u)=-\sum_{v\ \mathrm{internal}}
       \log(t_v^2-z_{v,0}^2-z_{v,1}^2).
\]
First leave all $D$ leaves free, with only the root fixed to one.
Put $J=\log_2D$.  At $u=0$, the analytic center of this full tree has
a common value $t_d^\circ$ at each depth $d$, where
\begin{equation}\label{eq:tree-center}
 (t_d^\circ)^2=\frac{D-2^d}{2^d(D-1)},\qquad0\leq d<J.
\end{equation}
To verify this, set $y_d=2^dt_d^2$, $y_0=1$, and $y_J=0$.
The barrier on depth-symmetric points, up to a constant, is
$-\sum_{d=0}^{J-1}2^d\log(y_d-y_{d+1})$.
Its minimizing gaps are $2^d/(D-1)$, giving
\eqref{eq:tree-center}.  This point is stationary in all tree variables:
automorphisms are transitive at each depth, so vanishing of each depth
derivative makes every coordinate derivative vanish.
Strict convexity of the sum of Lorentz barriers on the affine slice
then proves it is the unique center; each nonzero variable direction
affects at least one strictly convex cone term.  All leaf derivatives
vanish at $u=0$. This full-tree center lies on the affine slice fixing
the padded leaves to zero, so restriction preserves its minimizing
property and uniqueness. This argument does not require the padded
problem itself to retain the full tree symmetry.

Each leaf enters quadratically only in its bottom cone.  Hence at this
center
\[
 \nabla^2_{ut}\Psi=0,\qquad
 \nabla^2_{uu}\Psi=c_DI,\qquad c_D=2(D-1).
\]
After $u=Mx$, the $x$ Hessian block is $c_DM^TM$ with zero cross block.
For the objective $2e^Tu$, the first Newton direction in $x$ is
$(2\eta/c_D)M^{-1}e$.  For the normalized tilt, the squared decrement is
\begin{equation}\label{eq:tree-decrement}
 \Lambda_\eta^2=\frac{\eta^2}{c_D}
        \left(\frac54+\frac hG\Phi\right).
\end{equation}
The tree-variable block is public, has constant degree, and can be
preprocessed without hidden-input queries.  Thus full SQ access to the Newton
system is still simulable, but the gap between the low- and high-case
squared decrements is
$\Theta((\beta-\alpha)\eta^2\delta/(c_D\sqrt K))$.
The standard product barrier has parameter $2(D-1)$.
Neither a small condition number of the full lifted Hessian nor a
polylogarithmic iteration bound follows from the reduced clock block.
The norm tree is the standard nested Lorentz-cone representation of a
Euclidean norm \citep{AlizadehGoldfarb2003}, not a new cone lift.

\subsection{Sampling outputs and the rare-event cost}\label{subsec:cyclic-sampling}

Let $\Pi$ be the projection onto the coordinates $(t,0^n)$ with
$t$ in the hold segment.  The normalized
inverse history $v=x_*/R$ satisfies
\begin{equation}\label{eq:history-projector}
 \norm{\Pi v}^2=p\Phi^2.
\end{equation}
If $0\leq\varepsilon_v<1$ and
$\norm{\widehat x-x_*}\leq\varepsilon_v R$, then
\[
 \frac{(\sqrt p|\Phi|-\varepsilon_v)_+^2}{(1+\varepsilon_v)^2}
 \leq\frac{\norm{\Pi\widehat x}^2}{\norm{\widehat x}^2}
 \leq\frac{(\sqrt p|\Phi|+\varepsilon_v)^2}{(1-\varepsilon_v)^2}.
\]
For a fixed source gap $\beta>\alpha$, choosing
$\varepsilon_v\leq c(\beta-\alpha)\sqrt p$ with a sufficiently small
constant makes the low- and high-case probabilities of sampling a
hold-segment coordinate differ by $c' p$, with both $O(p)$.  By the variance bound for a Bernoulli mean,
averaging $O(1/p)$ independent samples therefore decides the source
promise with constant error.  It follows that generating one sample needs
\begin{equation}\label{eq:rare-sampling-lower}
 \Omega\left(p\frac{q^{\ell/2}}{r\ell}\right)
\end{equation}
queries in the two-function case.  With the parameters
\eqref{eq:scalar-clock-parameters}, $p=\Theta(\delta^2)$, so this bound
is \eqref{eq:scalar-coordinate-lower} times $\Theta(\delta^2)$.
The argument also applies to mixtures of such laws whose components each
come from a vector satisfying the same perturbation bound.

Equation~\eqref{eq:rare-sampling-lower} assumes that the overall sample
law satisfies this guarantee.  It also holds if failures are flagged and
retrying has bounded expected cost, or if unflagged failures
(contamination) have probability at most $c'p$ for a sufficiently
small constant.  A constant unflagged failure probability on each
fresh sample call is not enough: it can overwhelm the vanishing plateau
probability.  A different output model allows one
randomized setup that produces a fixed good sampler with constant
probability, whose repeated draws all have a good sample law.  Taking a constant number of independent setups and a
majority vote of their source decisions gives the usual bounded-error
reduction, provided that setup succeeds with probability at least
$1/2+c_{\rm setup}$ for an absolute constant $c_{\rm setup}>0$, and each
decision conditional on a good setup errs with probability at most
$c_{\rm setup}/2$.  Each setup then decides correctly with probability bounded above one
half, even if bad setups always err, so setup failure need not vanish.
With $Q_{\rm setup}$ and $Q_{\rm draw}$ the query costs of one setup
and one draw, and $Q_{\rm source}$ the randomized query lower
bound for the source problem, the resulting cost bound is
$Q_{\rm setup}+O(1/p)Q_{\rm draw}=\Omega(Q_{\rm source})$.
It gives no lower bound on a draw when preprocessing is not counted.
In Theorem~\ref{thm:solution-sampling}, the polynomial approximation is
deterministic, so its setup cannot fail; the rejection-sampling
randomness only draws samples from a fixed approximate vector.
If the sampler is promised relative vector error $\varepsilon_v$,
choosing instead $p=\Theta(\varepsilon_v)$ gives the weaker but useful
form
\[
 \Omega\left(\frac{\varepsilon_v q^{\ell/2}}{r\ell}\right),\qquad
 \ell=\frac{\log(1/\varepsilon_v)}{6\alpha_K}+O(1),
\]
under the same contamination condition, for vector error sufficiently
small relative to the fixed source gap $\beta-\alpha$. The sharper perturbation
inequality above also allows $p=\Theta(\varepsilon_v^2)$ with a sufficiently
large fixed prefactor, at the cost of a factor $\varepsilon_v^2$ from
repetition.  The scalar estimation in Theorem~\ref{thm:scalar-coordinate}
uses ordinary bounded error, with no repetition factor or vanishing
failure probability.

At fixed sufficiently small vector error and fixed $q$, choosing
$p=\Theta(1)$ gives $T=\Theta(K)$ and $\log N=\Theta(K)$.
The fixed-$k$ construction then gives a classical lower bound of
$\widetilde\Omega_k(N^{1-1/k})$ for sampling, with
$\kappa(M)=\Theta(\log N)$ and
$\kappa(H_0)=\Theta(\log^2N)$.  These condition numbers match the scales
shown to be necessary after Theorem~\ref{thm:solution-sampling}.
For full SQ access to the first Newton system, use $(H_0,M^Te)$; its
solution is $M^{-1}e/2$, and its full-SQ simulation was proved above.
An exactly feasible point with SOCP objective gap $\tau$ has relative
vector error at most $K_{\rm eff}\sqrt\tau\leq\sqrt{K\tau}$.
Thus, if $\tau\leq\varepsilon_v^2/K$, samples from such a point satisfy
this relative-error guarantee.

The family also has explicit upper bounds.  Classically,
reading all hidden sign tables costs $O_k(q^\ell)$ queries; computing
their circuit prefixes and sampling the public geometric clock gives
the exact inverse-state distribution in
$q^\ell\poly_k(r,\ell)$ arithmetic operations.  A designated matrix value query
recovers each sign, so this upper bound uses the same input interface.
Consequently the logarithm of the classical query complexity is
$\Theta(K\log s\log(1/\varepsilon_v))$ whenever its leading term dominates the
displayed repetition and denominator factors.
A quantum algorithm prepares the public clock state with geometric
amplitudes and applies each gate controlled on the clock having passed
that gate's position.  Forward and reverse segments
use each hidden sign gate twice in total. This prepares the exact
normalized history in the ideal circuit model. For any requested Euclidean
state error $0<\eta<1$, synthesize the public circuit to operator error at
most $\eta$: allocating error at most $\eta/G_{\rm public}$ to each of the
$G_{\rm public}$ public gates suffices by telescoping. The
implementation retains $O(k)$ hidden-input queries and public gate cost
polynomial in $n,T,\log(1/\eta)$. Taking $\eta\leq\varepsilon_v$
identifies the synthesized normalized state with a vector $\widehat x/R$
satisfying the relative-vector sampling tolerance above.
This algorithm uses the source-equivalent coherent oracle; a classical
SQ data structure alone does not provide its controlled sign queries.

A generic upper bound also holds for \emph{sampling} from an
approximately optimal feasible point.  The normal-equation polynomial
used in Theorem~\ref{thm:solution-sampling} can produce $z=P e$ with $\norm{Mz-e}\leq\xi$, at degree
$O(K\log(1/\xi))$.  This follows directly from the singular-value
residual bound for $M P$, not just the relative error in $z$.
For $\xi<1/2$, the normalized point
\[
 x=z/\norm{Mz},\qquad u=Mz/\norm{Mz}
\]
is exactly feasible for \eqref{eq:one-cone-program} and has objective
gap $O(\xi^2)$: its squared angular distance from $e$ is bounded by
a constant times $\norm{Mz-e}^2$.  Scaling $z$ does not change its
squared-coordinate sample law.  Setting $\xi=\Theta(\varepsilon_v/\sqrt K)$
therefore gives a sample from a point of gap $O(\varepsilon_v^2/K)$ with
exponential cost
$\exp(O(K\log(s+1)\log(\sqrt K/\varepsilon_v)))$, up to polynomial factors.
This bound covers samples from the $x$ block only; it does not provide
the normalization or the explicit coordinates of the feasible point.

\subsection{Relation to prior work}

Forrelation decision and its near-linear classical lower bounds are due
to Aaronson--Ambainis and Bansal--Sinha
\citep{AaronsonAmbainis2018,BansalSinha2021}.  The history-state sampling
obstruction is already present in Gr\o nlund--Larsen
\citep{GronlundLarsen2025}.  Ellipsoids, norm trees, and their SOCP
formulations are standard.  Apers and Gribling
\citep[Lemma 8.5]{ApersGribling2026} already prove a full-SQ classical
lower bound for constant-error LP optimal-value estimation.  Therefore
neither an SQ lower bound for optimization nor the hardness of scalar
optimization outputs is new here.  To the best of our knowledge, the
following are new: the joint dependence on conditioning, sparsity, and
accuracy for these specific formulations; the observation that
$G=\norm{M^{-T}a}$ is public, which permits the larger gap between
optimal values, and the optimality of this readout on the
plateau; and exact realizations of the hardness as a Newton decrement under
direct full SQ access to the Newton system and as the value of an
affine-slice LP with a parameter-one barrier, each under its stated
access model.

The coherent upper bounds are explicit algorithms for these structured
instances, not a general quantum interior-point method for
arbitrary SOCPs or LPs.  Quantum SOCP methods that reconstruct classical
Newton directions produce a different output
\citep{KerenidisPrakashSzilagyi2021}.  Building an explicit matrix table
or outputting an entire optimizer has dimension-dependent cost and lies
outside the scalar query models studied here.
